\section{递推建模与递归树}
\label{sec:04-Discrete-Mathematics/02-Recurrences/01}

给你一个数列，只告诉你后一项比前一项大 \(n\)，不告诉你第一项是多少。你能说出第 100 项的值吗？不能。递推规则只描述相邻之间的关系，它不选具体是哪条轨道——同一套生成规则下，不同初始条件跑出完全不同的数列。

这就是递推式的双面性格：局部规则+初始条件=唯一序列。缺任何一个，答案都不确定。

\subsection{递推式必须和初始条件一起出现}

拿最简单的递推 \(a_n=a_{n-1}+n\) 说事。如果你只知道这条规则，你可以写出：

\[
a_1=a_0+1,\quad a_2=a_1+2=a_0+1+2,\quad a_3=a_2+3=a_0+1+2+3.
\]

\(a_0\) 是一个你还没有填的数字。无论你取 \(a_0=0\) 还是 \(a_0=100\)，递推规则都照常工作——它只是忠实地在你填的起点上逐层累加。一旦你说 \(a_0=0\)，序列就被锁死了：\(a_n=0+1+2+\cdots+n=n(n+1)/2\)。

反过来说，如果你只拿到``第 5 项是 15''这个条件，没有递推规则，你也无法确定第 6 项——它可能接着加 6，也可能突然变成 0。规则和初始条件各自承担不同的责任：规则管``怎么变''，初始条件管``从哪开始''。

一个 \(k\) 阶递推（当前项依赖前 \(k\) 项）需要 \(k\) 个初值才能唯一确定整个序列。一阶只要一个起点，二阶要两个，以此类推。

建立递推模型的第一步是明确规模参数和最后一步。设 \(a_n\) 是长度为 \(n\) 的二进制串的数量。盯着最后一位看：删掉它之后剩下一个任意的长度 \(n-1\) 的二进制串，有 \(a_{n-1}\) 种；而被删掉的末位可以是 0 或 1，两种取法。于是 \(a_n=2a_{n-1}\)。再加初值 \(a_0=1\)（空串算 1 种），展开即得 \(a_n=2^n\)。递推在这个例子里只是``每次多一位，可能性翻倍''的形式化。

\subsection{分类最后一步产生组合递推}

再看一个不那么显然的例子。设 \(a_n\) 是不含连续两个 1 的长度 \(n\) 二进制串的数量。

写出小 \(n\) 来建立感觉：\(n=1\) 时，串 0 和串 1 都合法，\(a_1=2\)。\(n=2\) 时，合法串是 00, 01, 10——11 被禁了，\(a_2=3\)。\(n=3\) 时，合法串有 000, 001, 010, 100, 101，共 5 个。数列 \(2,3,5,\ldots\) 看起来有点眼熟。

现在建立递推。盯着长度为 \(n\) 的合法串的最后一位：

\begin{itemize}
\tightlist
\item
  末位为 0：前面 \(n-1\) 位可以是任何合法串，有 \(a_{n-1}\) 种。
\item
  末位为 1：倒数第二位绝对不能是 1（否则出现连续的 11），所以倒数第二位必须是 0。前面 \(n-2\) 位可以是任何合法串，有 \(a_{n-2}\) 种。
\end{itemize}

这两种情况互斥（末位不可能同时是 0 和 1），而且穷尽了所有可能。所以

\[
a_n=a_{n-1}+a_{n-2},
\]

配合初值 \(a_0=1\)（空串）、\(a_1=2\)。递推来自互斥且穷尽的结构分类，不是因为你认出了 Fibonacci 形状然后往上套公式。

\(a_0=1\) 可能让你不太舒服——长度为 0 的串只有一个，就是空串，空串显然不含连续两个 1。这个基例让递推能在 \(n=2\) 处正常工作：\(a_2=a_1+a_0=2+1=3\)，对上了。

\subsection{递归树看每层工作量}

递推不仅用来建立组合模型，更常见的用途是分析分治算法的运行时间。典型形式是

\[
T(n)=a\,T\!\left(\frac nb\right)+f(n).
\]

解释一下这个结构：一个规模 \(n\) 的问题被拆成 \(a\) 个子问题，每个规模 \(n/b\)；拆分和合并的额外工作量为 \(f(n)\)。

想知道 \(T(n)\) 大概是多少，一个直接的办法是把递推一层层展开，画成一棵树。根节点是规模 \(n\) 的原始问题，它做 \(f(n)\) 的工作，然后分出 \(a\) 个子节点。每个子节点规模 \(n/b\)，各做 \(f(n/b)\) 的工作，再分出 \(a\) 个孙节点。到第 \(j\) 层，一共有 \(a^j\) 个节点，每个节点规模 \(n/b^j\)，各自承担 \(f(n/b^j)\) 的工作。整层的工作量是

\[
a^j\,f\!\left(\frac{n}{b^j}\right).
\]

一直展开到叶节点——规模缩小到常数、不再继续拆分时停止。树的总工作量是各层工作量之和。

拿归并排序当实例。递推式 \(T(n)=2T(n/2)+cn\)，其中 \(cn\) 是拆分+合并的线性开销。这里 \(a=b=2\)，\(f(n)=cn\)。第 \(j\) 层有 \(2^j\) 个节点，每个规模 \(n/2^j\)，每节点开销 \(c\cdot n/2^j\)。该层总开销：

\[
2^j\cdot c\cdot\frac{n}{2^j}=cn.
\]

每一层的工作量都是 \(cn\)——不变。树有多少层？每次规模减半，从 \(n\) 缩到常数需要 \(\log_2 n\) 层。因此

\[
T(n)=cn\cdot\log_2 n+\text{叶节点成本},
\]

即 \(T(n)=\Theta(n\log n)\)。这里 \(\Theta\) 的意思是``同阶''：\(T(n)\) 与 \(n\log n\) 的比值被夹在两个正常数之间。

递归树的价值不只是算出答案，而是把总成本拆成两个可独立检查的来源：每层做多少工作、共有多少层。一旦看清楚这两项，你就能判断如果把问题拆得更碎（增大 \(b\)）或者拆出更多子问题（增大 \(a\)），总工作量会怎么变。这是在分析算法，而不是在套公式。

\subsection{猜出公式后仍需验证}

展开几项、看出规律、猜一个闭式——这个流程很自然，但猜完不能收工。归纳法负责把猜出的公式从``有道理''升级为``被证明''。

如果递推给的是不等式（例如 \(T(n)\le 2T(n/2)+cn\)），代入归纳假设时保留常数、检查基例是否被覆盖、确认不等号方向在代换中没有翻转——这些步骤一个都不能跳过。只比较最高阶项而忽略小规模处的常数，可能导致归纳步闭合失败：你在纸面上完成了一个看似自洽的推导，但它并没有真正被证明。

若 \(n\) 不是整齐的 \(b\) 的次幂（比如 \(n=10\)，\(b=3\)），取整带来的影响通常只涉及常数因子或低阶项，但需要说明，不能假装 \(n/b\) 总是整数。
